%%%%%%%%%%%%%%%%%%%%%%%%%%%%%%%%%%%%%%%%%%%%%%%%%%%%%%%%%%%%%%%%%%%%%%%%%%%%%%%%
\chapter{Yang--Mills Gauge Theory in the SHP Formalism}
\label{ch:shpym}
%%%%%%%%%%%%%%%%%%%%%%%%%%%%%%%%%%%%%%%%%%%%%%%%%%%%%%%%%%%%%%%%%%%%%%%%%%%%%%%%
\section{Introduction}
The gauge principle developed in the previous chapter is formulated
within ordinary four-dimensional space-time. Physical fields depend only
upon the coordinates
\[x^\mu,\qquad\mu=0,1,2,3,\]
and gauge transformations are local functions of these variables.
The Stueckelberg--Horwitz--Piron (SHP) theory introduces an additional
evolution parameter
\[\tau,\]
which parametrizes the dynamical evolution of events. The wave function
therefore depends upon five independent variables,
\[\Psi=\Psi(x^\mu,\tau).\]
Consequently, the local gauge transformation must also depend upon the
evolution parameter,
\[\boxed{\Psi(x,\tau)\rightarrow U(x,\tau)\Psi(x,\tau).}\]
This seemingly modest extension has profound consequences.
In ordinary Yang--Mills theory the derivative operator contains only the
four space-time derivatives,
\[\partial_\mu.\]
In the SHP theory there exists an additional derivative,
\[\partial_\tau,\]
which must transform covariantly under local gauge transformations.
Gauge invariance therefore requires introducing an additional gauge
potential associated with translations in the evolution parameter.
Instead of four gauge fields,
\[A_\mu,\]
the theory now possesses five,
\[\boxed{A_\alpha,\qquad\alpha=0,1,2,3,5.}\]
The fifth component,
\[A_5,\]
plays a role analogous to the scalar potential in ordinary
electrodynamics but is associated with evolution in the invariant
parameter rather than in coordinate time.
The purpose of this chapter is to develop the corresponding
five-dimensional Yang--Mills theory systematically from the gauge
principle.
%%%%%%%%%%%%%%%%%%%%%%%%%%%%%%%%%%%%%%%%%%%%%%%%%%%%%%%%%%%%%%%%%%%%%%%%%%%%%%%%
\section{The Five-Dimensional Gauge Principle}
%%%%%%%%%%%%%%%%%%%%%%%%%%%%%%%%%%%%%%%%%%%%%%%%%%%%%%%%%%%%%%%%%%%%%%%%%%%%%%%%
The fundamental dynamical variable of the SHP theory is the wave function
\[\Psi=\Psi(x^\mu,\tau),\]
which depends on the four space-time coordinates together with the
invariant evolution parameter
\[\tau .\]
Let the wave function transform according to an arbitrary compact Lie
group \(G\),
\[\boxed{\Psi(x,\tau)\longrightarrow\Psi'(x,\tau)=U(x,\tau)\Psi(x,\tau),}\]
where
\[U(x,\tau)=\exp\!\left(i\alpha^a(x,\tau)T^a\right).\]
Since the transformation parameters depend upon both \(x^\mu\) and
\(\tau\), neither ordinary space-time derivatives nor the
\(\tau\)-derivative transform covariantly.
%%%%%%%%%%%%%%%%%%%%%%%%%%%%%%%%%%%%%%%%%%%%%%%%%%%%%%%%%%%%%%%%%%%%%%%%%%%%%%%%
\subsection{Failure of Ordinary Derivatives}
The space-time derivative transforms as
\[\partial_\mu\Psi'=(\partial_\mu U)\Psi+U\partial_\mu\Psi,\]
while the evolution derivative satisfies
\[\partial_\tau\Psi'=(\partial_\tau U)\Psi+U\partial_\tau\Psi.\]
Both expressions contain inhomogeneous terms proportional to derivatives
of the gauge transformation.
Consequently,
\[\partial_\mu\Psi,\qquad\partial_\tau\Psi,\]
cannot represent observable quantities.
%%%%%%%%%%%%%%%%%%%%%%%%%%%%%%%%%%%%%%%%%%%%%%%%%%%%%%%%%%%%%%%%%%%%%%%%%%%%%%%%
\subsection{Five-Dimensional Covariant Derivative}
Gauge covariance is restored by introducing the five-component gauge
connection
\[A_\alpha=A_\alpha^aT^a,\qquad\alpha=0,1,2,3,5,\]
and defining
\[\boxed{D_\alpha=\partial_\alpha+igA_\alpha .}\]
Explicitly,
\[D_\mu=\partial_\mu+igA_\mu,\]
and
\[\boxed{D_5=\partial_\tau+igA_5.}\]
The defining requirement is
\[\boxed{D'_\alpha\Psi'=UD_\alpha\Psi .}\]
Exactly as in ordinary Yang--Mills theory, this requirement uniquely
determines the transformation law of the gauge fields.
%%%%%%%%%%%%%%%%%%%%%%%%%%%%%%%%%%%%%%%%%%%%%%%%%%%%%%%%%%%%%%%%%%%%%%%%%%%%%%%%
\subsection{Transformation Law}
Proceeding exactly as in Chapter~4,
\[(D'_\alpha U-U D_\alpha)\Psi=0.\]
Since this equation must hold for arbitrary wave functions,
\[D'_\alpha U=U D_\alpha.\]
Therefore,
\[\boxed{A'_\alpha=UA_\alpha U^{-1}-\frac{i}{g}(\partial_\alpha U)U^{-1}.}\]
Remarkably,
this single equation simultaneously describes the transformation of the
ordinary four-vector potential
\[A_\mu,\]
and the new fifth component
\[A_5.\]
%%%%%%%%%%%%%%%%%%%%%%%%%%%%%%%%%%%%%%%%%%%%%%%%%%%%%%%%%%%%%%%%%%%%%%%%%%%%%%%%
\subsection{Abelian Limit}
For
\[G=U(1),\]
all commutators vanish,
and the transformations reduce to
\[A'_\mu=A_\mu-\partial_\mu\alpha,\]
together with
\[\boxed{A'_5=A_5-\partial_\tau\alpha.}\]
The second equation has no analogue in ordinary electrodynamics.
It is a direct consequence of allowing local gauge transformations to
depend upon the invariant evolution parameter.
%%%%%%%%%%%%%%%%%%%%%%%%%%%%%%%%%%%%%%%%%%%%%%%%%%%%%%%%%%%%%%%%%%%%%%%%%%%%%%%%
\subsection{Interpretation of the Fifth Gauge Field}
The field
\[A_5(x,\tau)\]
is not an additional spatial component of the gauge potential.
Instead,
it is associated with parallel transport in the direction of the
evolution parameter.
In four-dimensional gauge theory,
the connection specifies how internal states are compared at neighboring
space-time points.
In SHP theory,
comparison is also required between neighboring events having different
values of
\[\tau.\]
The field
\[A_5\]
provides precisely this additional connection.
It is therefore the gauge potential associated with evolution in the
world-time parameter rather than in physical space-time.
%%%%%%%%%%%%%%%%%%%%%%%%%%%%%%%%%%%%%%%%%%%%%%%%%%%%%%%%%%%%%%%%%%%%%%%%%%%%%%%%
\section{Five-Dimensional Gauge Curvature}
%%%%%%%%%%%%%%%%%%%%%%%%%%%%%%%%%%%%%%%%%%%%%%%%%%%%%%%%%%%%%%%%%%%%%%%%%%%%%%%%
The gauge connection introduced in the previous section defines a
covariant derivative
\[D_\alpha=\partial_\alpha+igA_\alpha,\qquad\alpha=0,1,2,3,5.\]
Exactly as in ordinary Yang--Mills theory, the gauge curvature is defined
through the commutator of covariant derivatives,
\[\boxed{[D_\alpha,D_\beta]=igF_{\alpha\beta}.}\]
Our task is to determine the explicit form of the tensor
\(F_{\alpha\beta}\).
%%%%%%%%%%%%%%%%%%%%%%%%%%%%%%%%%%%%%%%%%%%%%%%%%%%%%%%%%%%%%%%%%%%%%%%%%%%%%%%%
\subsection{Derivation from the Covariant Derivative}
Expanding the commutator,
\[[D_\alpha,D_\beta]=[\partial_\alpha+igA_\alpha, \partial_\beta+igA_\beta].\]
Using
\[[\partial_\alpha,\partial_\beta]=0,\]
one finds
\[[D_\alpha,D_\beta]=ig(\partial_\alpha A_\beta-\partial_\beta A_\alpha)
-g^2[A_\alpha,A_\beta].\]
Therefore,
\[\boxed{F_{\alpha\beta}=\partial_\alpha A_\beta
-\partial_\beta A_\alpha-ig[A_\alpha,A_\beta].}\]
This has precisely the same mathematical form as the ordinary
Yang--Mills curvature.
The difference lies entirely in the enlarged range of the indices.
%%%%%%%%%%%%%%%%%%%%%%%%%%%%%%%%%%%%%%%%%%%%%%%%%%%%%%%%%%%%%%%%%%%%%%%%%%%%%%%%
\subsection{Decomposition of the Curvature}
Since
\[\alpha=(\mu,5),\]
the curvature naturally decomposes into three sectors,
\[F_{\mu\nu},\qquad F_{\mu5},\qquad F_{55}.\]
Each possesses a distinct physical interpretation.
%%%%%%%%%%%%%%%%%%%%%%%%%%%%%%%%%%%%%%%%%%%%%%%%%%%%%%%%%%%%%%%%%%%%%%%%%%%%%%%%
\subsubsection{Space-Time Components}
The purely space-time components are
\[\boxed{F_{\mu\nu}=\partial_\mu A_\nu-\partial_\nu A_\mu-ig[A_\mu,A_\nu].}\]
These coincide exactly with the ordinary Yang--Mills field tensor.
Consequently,
all familiar electromagnetic or non-Abelian interactions are recovered
without modification.
%%%%%%%%%%%%%%%%%%%%%%%%%%%%%%%%%%%%%%%%%%%%%%%%%%%%%%%%%%%%%%%%%%%%%%%%%%%%%%%%
\subsubsection{Mixed Components}
The mixed components are
\[\boxed{F_{\mu5}=\partial_\mu A_5-\partial_\tau A_\mu-ig[A_\mu,A_5].}\]
This tensor has no analogue in ordinary four-dimensional gauge theory.
It contains three qualitatively different contributions.
The first,
\[\partial_\mu A_5,\]
describes the variation of the fifth gauge potential throughout
space-time.
The second,
\[-\partial_\tau A_\mu,\]
measures the explicit evolution of the ordinary gauge field with respect
to the invariant parameter.
The final commutator,
\[-ig[A_\mu,A_5],\]
couples the ordinary and fifth gauge potentials through the non-Abelian
Lie algebra.
%%%%%%%%%%%%%%%%%%%%%%%%%%%%%%%%%%%%%%%%%%%%%%%%%%%%%%%%%%%%%%%%%%%%%%%%%%%%%%%%
\subsubsection{Fifth Component}
Because the curvature tensor is antisymmetric,
\[F_{55}=0.\]
No independent fifth-fifth field strength exists.
%%%%%%%%%%%%%%%%%%%%%%%%%%%%%%%%%%%%%%%%%%%%%%%%%%%%%%%%%%%%%%%%%%%%%%%%%%%%%%%%
\subsection{Gauge Covariance}
The curvature transforms according to
\[\boxed{F_{\alpha\beta}\rightarrow UF_{\alpha\beta}U^{-1}.}\]
Consequently,
all gauge-invariant observables may be constructed from traces such as
\[\operatorname{Tr}(F_{\alpha\beta}F^{\alpha\beta}).\]
The extension from four to five dimensions therefore preserves the
entire gauge structure of Yang--Mills theory.
%%%%%%%%%%%%%%%%%%%%%%%%%%%%%%%%%%%%%%%%%%%%%%%%%%%%%%%%%%%%%%%%%%%%%%%%%%%%%%%%
\subsection{Abelian Limit}
For
\[G=U(1),\]
the commutators vanish.
The mixed field becomes
\[\boxed{F_{\mu5}=\partial_\mu A_5-\partial_\tau A_\mu.}\]
This expression is the SHP analogue of the electromagnetic field tensor.
Unlike Maxwell theory,
it contains derivatives with respect to the invariant evolution
parameter.
Consequently,
the Abelian SHP theory already predicts new dynamical effects absent from
ordinary electrodynamics.
%%%%%%%%%%%%%%%%%%%%%%%%%%%%%%%%%%%%%%%%%%%%%%%%%%%%%%%%%%%%%%%%%%%%%%%%%%%%%%%%
\subsection{Geometrical Interpretation}
The ordinary Yang--Mills curvature measures the failure of parallel
transport around infinitesimal loops lying entirely within space-time.
The tensor
\[F_{\mu5},\]
extends this concept.
It measures the failure of parallel transport around an infinitesimal
loop whose edges lie partly in space-time and partly along the evolution
parameter.
Thus,
the mixed curvature describes the coupling between local gauge rotations
generated by displacement in space-time and those generated by evolution
in world time.
It is therefore the genuinely new geometric object introduced by the SHP
gauge principle.
%%%%%%%%%%%%%%%%%%%%%%%%%%%%%%%%%%%%%%%%%%%%%%%%%%%%%%%%%%%%%%%%%%%%%%%%%%%%%%%%
\section{The Physical Interpretation of the Fifth Gauge Potential}
%%%%%%%%%%%%%%%%%%%%%%%%%%%%%%%%%%%%%%%%%%%%%%%%%%%%%%%%%%%%%%%%%%%%%%%%%%%%%%%%
The appearance of the fifth gauge potential
\[A_5(x,\tau)\]
is the most distinctive feature of the SHP gauge theory.
Mathematically, its existence follows immediately from the requirement
of local gauge invariance in the extended space
\[(x^\mu,\tau).\]
Physically, however, its interpretation is considerably more subtle.
Unlike the four-vector potential,
\[A_\mu,\]
which is associated with parallel transport in space-time,
the field
\[A_5\]
governs parallel transport in the direction of the invariant evolution
parameter.
Its appearance therefore reflects the dynamical structure of SHP theory
rather than the geometry of ordinary space-time.
%%%%%%%%%%%%%%%%%%%%%%%%%%%%%%%%%%%%%%%%%%%%%%%%%%%%%%%%%%%%%%%%%%%%%%%%%%%%%%%%
\subsection{Appearance in the Hamiltonian}
The gauge-invariant SHP Hamiltonian has the form
\[\boxed{K=\frac{1}{2M}(p_\mu-gA_\mu)(p^\mu-gA^\mu)-gA_5.}\]
The first term represents the kinetic energy modified by the ordinary
gauge interaction.
The second term is qualitatively different.
Unlike
\[A_\mu,\]
which enters through the generalized momentum,
the field
\[A_5\]
appears directly in the Hamiltonian itself.
Consequently,
it contributes directly to the generator of evolution with respect to the
world-time parameter
\[\tau.\]
%%%%%%%%%%%%%%%%%%%%%%%%%%%%%%%%%%%%%%%%%%%%%%%%%%%%%%%%%%%%%%%%%%%%%%%%%%%%%%%%
\subsection{Dimensional Analysis}
Since the Hamiltonian satisfies
\[i\frac{\partial\Psi}{\partial\tau}=K\Psi,\]
the quantity
\[K\]
has the dimensions conjugate to the evolution parameter.
Therefore
\[gA_5\]
must possess exactly the same dimensions.
This differs fundamentally from the ordinary gauge potential,
whose dimensions are fixed through its coupling to the four-momentum,
\[p_\mu-gA_\mu.\]
The fifth gauge field therefore belongs to a different dynamical sector.
%%%%%%%%%%%%%%%%%%%%%%%%%%%%%%%%%%%%%%%%%%%%%%%%%%%%%%%%%%%%%%%%%%%%%%%%%%%%%%%%
\subsection{Generator of World-Time Evolution}
The ordinary gauge field modifies the particle trajectory through the
Lorentz force,
\[M\ddot x^\mu=gF^{\mu\nu}\dot x_\nu.\]
The field
\[A_5,\]
however,
changes the Hamiltonian itself.
Consequently,
it affects every dynamical equation obtained from Hamilton's principle.
Rather than producing an additional force in space-time,
it changes the evolution generated by
\[K.\]
%%%%%%%%%%%%%%%%%%%%%%%%%%%%%%%%%%%%%%%%%%%%%%%%%%%%%%%%%%%%%%%%%%%%%%%%%%%%%%%%
\subsection{Off-Shell Dynamics}
One of the defining characteristics of SHP theory is that the invariant
mass
\[m^2=p^\mu p_\mu\]
is not constrained to remain constant.
Instead,
the particle evolves off shell.
Since
\[A_5\]
enters directly into the Hamiltonian,
it contributes naturally to the evolution of the invariant mass.
Thus,
the fifth gauge potential provides a mechanism through which gauge
interactions may exchange invariant mass with the particle while
preserving the gauge symmetry of the theory.
%%%%%%%%%%%%%%%%%%%%%%%%%%%%%%%%%%%%%%%%%%%%%%%%%%%%%%%%%%%%%%%%%%%%%%%%%%%%%%%%
\subsection{Gauge Interpretation}
The gauge transformation
\[A_5'=UA_5U^{-1}-\frac{i}{g}(\partial_\tau U)U^{-1}\]
shows that
\[A_5\]
is the connection associated with infinitesimal displacements in the
world-time direction.
It therefore plays precisely the same geometric role along
\[\tau\]
as
\[A_\mu\]
plays along the four space-time directions.
The complete gauge connection may therefore be written as
\[A_\alpha=(A_\mu,A_5),\]
with no distinction at the level of the underlying geometry.
%%%%%%%%%%%%%%%%%%%%%%%%%%%%%%%%%%%%%%%%%%%%%%%%%%%%%%%%%%%%%%%%%%%%%%%%%%%%%%%%
\subsection{Comparison with Scalar Fields}
The fifth gauge potential should not be confused with an ordinary scalar
field.
A scalar field transforms according to
\[\phi\rightarrow U\phi.\]
The field
\[A_5\]
instead transforms as
\[A_5\rightarrow UA_5U^{-1}-\frac{i}{g}(\partial_\tau U)U^{-1}.\]
The inhomogeneous term is the characteristic signature of a gauge
connection.Therefore,although
\[A_5\]
is a Lorentz scalar,
it is not a scalar field in the ordinary sense.
It is the fifth component of a gauge connection.
%%%%%%%%%%%%%%%%%%%%%%%%%%%%%%%%%%%%%%%%%%%%%%%%%%%%%%%%%%%%%%%%%%%%%%%%%%%%%%%%
\subsection{Comparison with Kaluza--Klein Theory}
The additional coordinate appearing in SHP theory should not be confused
with the compact fifth dimension of Kaluza--Klein models.
In Kaluza--Klein theory,
the extra coordinate is interpreted as an additional geometric dimension
of space-time.
In SHP theory,
the parameter
\[\tau\]
is an invariant evolution parameter rather than a physical coordinate.
Consequently,the field
\[A_5\]
is associated with evolution in world time,
not with motion through an additional spatial dimension.
Although the mathematical structures are superficially similar,
their physical interpretations are fundamentally different.
%%%%%%%%%%%%%%%%%%%%%%%%%%%%%%%%%%%%%%%%%%%%%%%%%%%%%%%%%%%%%%%%%%%%%%%%%%%%%%%%
\section{Gauge-Covariant Hamilton Equations}
%%%%%%%%%%%%%%%%%%%%%%%%%%%%%%%%%%%%%%%%%%%%%%%%%%%%%%%%%%%%%%%%%%%%%%%%%%%%%%%%
The gauge-invariant Hamiltonian for a classical particle interacting
with a non-Abelian SHP gauge field is
\[\boxed{K=\frac{1}{2M}\left(p_\mu-gA_\mu^aQ^a\right)\left(p^\mu
-gA^{\mu b}Q^b\right)-gA_5^aQ^a.}\]
For convenience we introduce the gauge-covariant momentum
\[\boxed{\Pi_\mu=p_\mu-gA_\mu^aQ^a.}\]
The Hamiltonian becomes
\[\boxed{K=\frac{1}{2M}\Pi^\mu\Pi_\mu-gA_5^aQ^a.}\]
%%%%%%%%%%%%%%%%%%%%%%%%%%%%%%%%%%%%%%%%%%%%%%%%%%%%%%%%%%%%%%%%%%%%%%%%%%%%%%%%
\subsection{Equation for the Coordinates}
Hamilton's first equation is
\[\dot x^\mu=\frac{\partial K}{\partial p_\mu}.\]
Since
\[\frac{\partial\Pi_\nu}{\partial p_\mu}=\delta^\mu_\nu,\]
we obtain
\[\boxed{\dot x^\mu=\frac{\Pi^\mu}{M}.}\]
Thus,
exactly as in ordinary Yang--Mills theory,
the covariant momentum determines the physical four-velocity,
\[\boxed{\Pi^\mu=M\dot x^\mu.}\]
%%%%%%%%%%%%%%%%%%%%%%%%%%%%%%%%%%%%%%%%%%%%%%%%%%%%%%%%%%%%%%%%%%%%%%%%%%%%%%%%
\subsection{Equation for the Canonical Momentum}
The canonical momentum satisfies
\[\dot p_\mu=-\frac{\partial K}{\partial x^\mu}.\]
Since both
\[A_\mu^a\]
and
\[A_5^a\]
depend upon
\[x^\mu,\]
one finds
\[\boxed{\dot p_\mu=\frac{g}{M}\Pi^\nu Q^a\partial_\mu A_\nu^a
+gQ^a\partial_\mu A_5^a.}\]
Already at this stage we observe a new contribution absent from ordinary
Yang--Mills theory.
The fifth gauge potential contributes directly to the force.
%%%%%%%%%%%%%%%%%%%%%%%%%%%%%%%%%%%%%%%%%%%%%%%%%%%%%%%%%%%%%%%%%%%%%%%%%%%%%%%%
\subsection{Evolution of the Covariant Momentum}
Differentiating
\[\Pi_\mu=p_\mu-gA_\mu^aQ^a,\]
gives
\[\begin{aligned}
\dot\Pi_\mu=&\dot p_\mu-g\dot x^\nu\partial_\nu A_\mu^aQ^a\\
&-gA_\mu^a\dot Q^a-g\partial_\tau A_\mu^aQ^a.
\end{aligned}\]
The final term appears because
\[A_\mu=A_\mu(x,\tau),
\]
depends explicitly upon the evolution parameter.
Substituting
Hamilton's equation for
\[\dot p_\mu\]
and the Wong equation for
\[\dot Q^a,\]
all terms reorganize into the five-dimensional field tensor.
One finally obtains
\[\boxed{\dot\Pi_\mu=gQ^aF_{\mu\alpha}^a\dot x^\alpha.}\]
Here
\[\alpha=0,1,2,3,5,\]
and
\[\dot x^5\equiv 1,\]
because
\[x^5=\tau.\]
%%%%%%%%%%%%%%%%%%%%%%%%%%%%%%%%%%%%%%%%%%%%%%%%%%%%%%%%%%%%%%%%%%%%%%%%%%%%%%%%
\subsection{Generalized Lorentz Force}
Since
\[\Pi^\mu=M\dot x^\mu,\]
the equation of motion becomes
\[\boxed{M\ddot x^\mu=gQ^aF^{a\mu}{}_{\alpha}\dot x^\alpha.}\]
Separating the index,
\[\alpha=(\nu,5),\]
gives
\[\boxed{M\ddot x^\mu=gQ^aF^{a\mu}{}_{\nu}\dot x^\nu+gQ^aF^{a\mu5}.}\]
The first term is the ordinary Wong force.
The second,
proportional to
\[F^{\mu5},\]
is entirely new.
It is the characteristic SHP correction.
%%%%%%%%%%%%%%%%%%%%%%%%%%%%%%%%%%%%%%%%%%%%%%%%%%%%%%%%%%%%%%%%%%%%%%%%%%%%%%%%
\subsection{Interpretation}
The tensor
\[F_{\mu5}\]
acts as an additional gauge force generated by variations of the gauge
connection with respect to the invariant evolution parameter.
Consequently,
the particle dynamics is governed not only by the ordinary gauge field,
but also by the evolution of that field in world time.
The generalized Lorentz force therefore naturally separates into two
parts,
\[\boxed{\text{ordinary Yang--Mills force}+\text{world-time force}.}\]
The latter has no analogue in conventional gauge theory.
%%%%%%%%%%%%%%%%%%%%%%%%%%%%%%%%%%%%%%%%%%%%%%%%%%%%%%%%%%%%%%%%%%%%%%%%%%%%%%%%
\section{Analysis of the Generalized Lorentz Force}
%%%%%%%%%%%%%%%%%%%%%%%%%%%%%%%%%%%%%%%%%%%%%%%%%%%%%%%%%%%%%%%%%%%%%%%%%%%%%%%%
The equation of motion derived in the previous section,
\[M\ddot x^\mu=gQ^aF^{a\mu}{}_{\nu}\dot x^\nu+gQ^aF^{a\mu5},\]
contains the ordinary Wong force together with an additional
contribution arising from the mixed components of the five-dimensional
field tensor.
The purpose of this section is to examine the physical content of this
new interaction.
%%%%%%%%%%%%%%%%%%%%%%%%%%%%%%%%%%%%%%%%%%%%%%%%%%%%%%%%%%%%%%%%%%%%%%%%%%%%%%%%
\subsection{Decomposition of the Mixed Field}
Using the definition of the five-dimensional curvature,
\[F_{\mu5}^a=\partial_\mu A_5^a-\partial_\tau A_\mu^a-gf^{abc}A_\mu^bA_5^c,\]
the generalized force becomes
\[M\ddot x^\mu=gQ^aF^{a\mu}{}_{\nu}\dot x^\nu+gQ^a
\left(\partial^\mu A_5^a-\partial_\tau A^{a\mu}-gf^{abc}A^{b\mu}A_5^c\right).\]
The SHP correction therefore consists of three qualitatively distinct
contributions.
%%%%%%%%%%%%%%%%%%%%%%%%%%%%%%%%%%%%%%%%%%%%%%%%%%%%%%%%%%%%%%%%%%%%%%%%%%%%%%%%
\subsection{Gradient of the Fifth Potential}
The first term,
\[gQ^a\partial^\mu A_5^a,\]
has the structure of the gradient of a potential.
Unlike the ordinary gauge interaction,
which acts through the covariant momentum,
this contribution originates directly from the appearance of
\[A_5\]
in the Hamiltonian.
It represents the force associated with spatial variations of the
generator of evolution.
%%%%%%%%%%%%%%%%%%%%%%%%%%%%%%%%%%%%%%%%%%%%%%%%%%%%%%%%%%%%%%%%%%%%%%%%%%%%%%%%
\subsection{Evolution of the Gauge Field}
The second contribution,
\[-gQ^a\partial_\tau A^{a\mu},\]
is entirely new.
It appears because the gauge field itself evolves with respect to the
world-time parameter.
If
\[\partial_\tau A_\mu^a=0,\]
the theory reduces locally to ordinary Yang--Mills dynamics.
Whenever
\[\partial_\tau A_\mu^a\neq0,\]
the particle experiences an additional force arising solely from the
evolution of the gauge field.
%%%%%%%%%%%%%%%%%%%%%%%%%%%%%%%%%%%%%%%%%%%%%%%%%%%%%%%%%%%%%%%%%%%%%%%%%%%%%%%%
\subsection{Non-Abelian Mixing}
The third contribution,
\[-g^2f^{abc}Q^aA^{b\mu}A_5^c,\]
exists only for non-Abelian gauge groups.
It couples the ordinary gauge potential to the fifth gauge potential.
This interaction disappears completely in the Abelian limit,
\[f^{abc}=0.\]
Thus it represents a genuinely non-Abelian SHP effect.
%%%%%%%%%%%%%%%%%%%%%%%%%%%%%%%%%%%%%%%%%%%%%%%%%%%%%%%%%%%%%%%%%%%%%%%%%%%%%%%%
\subsection{Static Limit}
Suppose
\[\partial_\tau A_\mu^a=0.\]
Then
\[F_{\mu5}^a=\partial_\mu A_5^a-gf^{abc}A_\mu^bA_5^c.\]
If, in addition,
\[A_5^a=0,\]
the mixed field vanishes,
\[F_{\mu5}=0,\]
and the generalized Lorentz force reduces exactly to the classical Wong
equation.
Ordinary Yang--Mills theory is therefore recovered as a special case of
the SHP dynamics.
%%%%%%%%%%%%%%%%%%%%%%%%%%%%%%%%%%%%%%%%%%%%%%%%%%%%%%%%%%%%%%%%%%%%%%%%%%%%%%%%
\subsection{Pure Fifth-Potential Interaction}
Consider the opposite limit,
\[A_\mu^a=0,\]
while
\[A_5^a\neq0.\]
Then
\[F_{\mu5}^a=\partial_\mu A_5^a,\]
and
\[M\ddot x^\mu=gQ^a\partial^\mu A_5^a.\]
The particle experiences a force even though no ordinary gauge field is
present.
This interaction is entirely absent from conventional gauge theory.
%%%%%%%%%%%%%%%%%%%%%%%%%%%%%%%%%%%%%%%%%%%%%%%%%%%%%%%%%%%%%%%%%%%%%%%%%%%%%%%%
\subsection{Slow Evolution Approximation}
If the gauge field varies slowly with world time,
\[\left|\partial_\tau A_\mu\right|\ll\left|\partial_\nu A_\mu\right|,\]
the SHP correction may be regarded as a perturbation of the ordinary
Yang--Mills dynamics.
This approximation is expected to be useful whenever the evolution in
\(\tau\) occurs on scales much longer than the characteristic
space-time scales of the interaction.
%%%%%%%%%%%%%%%%%%%%%%%%%%%%%%%%%%%%%%%%%%%%%%%%%%%%%%%%%%%%%%%%%%%%%%%%%%%%%%%%
\subsection{Abelian Limit}
For
\[G=U(1),\]
the structure constants vanish,
\[f^{abc}=0,\]
and the generalized force becomes
\[M\ddot x^\mu=eF^{\mu\nu}\dot x_\nu+e\left(\partial^\mu A_5
-\partial_\tau A^\mu\right).\]
Thus even the Abelian SHP theory predicts new interactions generated by
the fifth gauge field.
%%%%%%%%%%%%%%%%%%%%%%%%%%%%%%%%%%%%%%%%%%%%%%%%%%%%%%%%%%%%%%%%%%%%%%%%%%%%%%%%
\subsection{Physical Interpretation}
The generalized Lorentz force naturally separates into two distinct
parts.
The first,
\[gQ^aF^{a\mu}{}_{\nu}\dot x^\nu,\]
describes the interaction of the particle with the ordinary gauge
curvature in space-time.
The second,
\[gQ^aF^{a\mu5},\]
describes the interaction with the evolution of the gauge connection in
the world-time direction.
The complete dynamics therefore depends upon the curvature of the gauge
connection in the extended manifold
\[M_4\times\mathbb{R}_\tau,\]
rather than solely upon the curvature in space-time.
%%%%%%%%%%%%%%%%%%%%%%%%%%%%%%%%%%%%%%%%%%%%%%%%%%%%%%%%%%%%%%%%%%%%%%%%%%%%%%%%
\section{Evolution of the Invariant Mass}
%%%%%%%%%%%%%%%%%%%%%%%%%%%%%%%%%%%%%%%%%%%%%%%%%%%%%%%%%%%%%%%%%%%%%%%%%%%%%%%%
One of the defining characteristics of the SHP theory is that the
particle is not constrained to remain on a fixed mass shell.
The invariant
\[p^\mu p_\mu\]
is therefore no longer a constant of motion.
Instead,
its evolution is determined dynamically by the Hamilton equations.
%%%%%%%%%%%%%%%%%%%%%%%%%%%%%%%%%%%%%%%%%%%%%%%%%%%%%%%%%%%%%%%%%%%%%%%%%%%%%%%%
\subsection{Evolution of the Canonical Momentum}
The generalized Hamilton equation is
\[\dot p_\mu=gQ^a\partial_\mu A_5^a+\frac{g}{M}Q^a\Pi^\nu\partial_\mu A_\nu^a.\]
Consequently,
\[\frac{d}{d\tau}(p^\mu p_\mu)=2p^\mu\dot p_\mu.\]
Substituting the Hamilton equation gives
\[\boxed{\frac{d}{d\tau}(p^2)=2gp^\mu Q^a\partial_\mu A_5^a
+\frac{2g}{M}p^\mu Q^a\Pi^\nu\partial_\mu A_\nu^a.}\]
This equation shows explicitly that the invariant mass changes whenever
the gauge potentials possess nontrivial space-time dependence.
%%%%%%%%%%%%%%%%%%%%%%%%%%%%%%%%%%%%%%%%%%%%%%%%%%%%%%%%%%%%%%%%%%%%%%%%%%%%%%%%
\subsection{Contribution of the Fifth Potential}
The first term,
\[2gp^\mu Q^a\partial_\mu A_5^a,\]
originates entirely from the fifth gauge potential.
It has no analogue in ordinary Yang--Mills theory.
Spatial variations of
\[A_5\]
therefore produce evolution of the invariant mass.
%%%%%%%%%%%%%%%%%%%%%%%%%%%%%%%%%%%%%%%%%%%%%%%%%%%%%%%%%%%%%%%%%%%%%%%%%%%%%%%%
\subsection{Gauge-Field Contribution}
The second term,
\[\frac{2g}{M}p^\mu Q^a\Pi^\nu\partial_\mu A_\nu^a,\]
describes the influence of the ordinary gauge field.
Even when
\[A_5=0,\]
space-time variations of the gauge field can modify the invariant mass
through the off-shell dynamics.
%%%%%%%%%%%%%%%%%%%%%%%%%%%%%%%%%%%%%%%%%%%%%%%%%%%%%%%%%%%%%%%%%%%%%%%%%%%%%%%%
\subsection{Ordinary Relativistic Mechanics}
If
\[A_\mu=A_5=0,\]
Hamilton's equations reduce to
\[\dot p_\mu=0,\]
and therefore
\[\boxed{\frac{d}{d\tau}(p^2)=0.}\]
The particle remains on a fixed mass shell.
Ordinary relativistic mechanics is recovered immediately.
%%%%%%%%%%%%%%%%%%%%%%%%%%%%%%%%%%%%%%%%%%%%%%%%%%%%%%%%%%%%%%%%%%%%%%%%%%%%%%%%
\subsection{Interpretation}
The evolution equation demonstrates that off-shell motion is not an
arbitrary feature of SHP mechanics.
Rather,
it arises naturally from the interaction of the particle with the
five-dimensional gauge connection.
The fifth gauge potential modifies not only the trajectory of the
particle,
but also the evolution of its invariant mass.
%%%%%%%%%%%%%%%%%%%%%%%%%%%%%%%%%%%%%%%%%%%%%%%%%%%%%%%%%%%%%%%%%%%%%%%%%%%%%%%%
\section{Canonical Structure of SHP Gauge Theory}
%%%%%%%%%%%%%%%%%%%%%%%%%%%%%%%%%%%%%%%%%%%%%%%%%%%%%%%%%%%%%%%%%%%%%%%%%%%%%%%%
The preceding sections established the gauge-covariant Hamiltonian
equations for a classical particle interacting with a five-dimensional
Yang--Mills field. We now examine the underlying canonical structure.
The central question is whether the enlarged gauge symmetry modifies the
symplectic geometry of the SHP phase space.
%%%%%%%%%%%%%%%%%%%%%%%%%%%%%%%%%%%%%%%%%%%%%%%%%%%%%%%%%%%%%%%%%%%%%%%%%%%%%%%%
\subsection{Extended Phase Space}
The dynamical variables of the particle are
\[(x^\mu,p_\mu,Q^a),\]
together with the evolution parameter
\[\tau.\]
The phase space is therefore
\[\Gamma=T^*M\times{\cal O},\]
where
\[{\cal O}\]
denotes the coadjoint orbit associated with the internal charge.
The evolution parameter is not itself a canonical coordinate.
Rather,it labels the Hamiltonian flow on the extended phase space.
%%%%%%%%%%%%%%%%%%%%%%%%%%%%%%%%%%%%%%%%%%%%%%%%%%%%%%%%%%%%%%%%%%%%%%%%%%%%%%%%
\subsection{Canonical One-Form}
The canonical one-form is
\[\boxed{\Theta=p_\mu dx^\mu+\Theta_{\rm color}-K\,d\tau .}\]
The first term is the ordinary Liouville one-form.
The second is the Kirillov--Kostant one-form on the coadjoint orbit.
The final term generates evolution with respect to the invariant
parameter.
The complete dynamics follows from the action
\[S=\int\Theta.\]
%%%%%%%%%%%%%%%%%%%%%%%%%%%%%%%%%%%%%%%%%%%%%%%%%%%%%%%%%%%%%%%%%%%%%%%%%%%%%%%%
\subsection{Gauge-Covariant Canonical Form}
Minimal coupling replaces
\[p_\mu\rightarrow\Pi_\mu=p_\mu-gA_\mu^aQ^a.\]
The canonical one-form therefore becomes
\[\boxed{\Theta=\Pi_\mu dx^\mu+gA_\mu^aQ^a dx^\mu+\Theta_{\rm color}-K\,d\tau .}\]
Using
\[K=\frac{\Pi^2}{2M}-gA_5^aQ^a,\]
the action may be written as
\[\boxed{S=\int\left(\Pi_\mu dx^\mu+\Theta_{\rm color}-\frac{\Pi^2}{2M}d\tau
+gA_\alpha^aQ^a dx^\alpha\right),}\]
where
\[dx^5=d\tau.\]
This expression treats all five components of the gauge connection on an
equal footing.
%%%%%%%%%%%%%%%%%%%%%%%%%%%%%%%%%%%%%%%%%%%%%%%%%%%%%%%%%%%%%%%%%%%%%%%%%%%%%%%%
\subsection{Symplectic Two-Form}
The symplectic structure follows from
\[\Omega=-d\Theta .\]
After differentiation,
the result separates naturally into three contributions,
\[\boxed{\Omega=\Omega_{\rm space}+\Omega_{\rm color}+\Omega_{\rm gauge}.}\]
The first term is the ordinary canonical symplectic form,
\[dp_\mu\wedge dx^\mu.\]
The second is the Kirillov--Kostant form on the color orbit.
The third is generated by the gauge interaction,
\[gQ^aF_{\alpha\beta}^adx^\alpha\wedge dx^\beta.\]
Thus the gauge curvature appears directly as part of the symplectic
geometry of the theory.
%%%%%%%%%%%%%%%%%%%%%%%%%%%%%%%%%%%%%%%%%%%%%%%%%%%%%%%%%%%%%%%%%%%%%%%%%%%%%%%%
\subsection{Hamiltonian Vector Field}
The Hamiltonian flow is determined by
\[i_{X_K}\Omega=dK.\]
The corresponding vector field reproduces
\[\dot x^\mu,\qquad\dot p_\mu,\qquad\dot Q^a,\]
and therefore generates simultaneously
the generalized Lorentz force,
the Wong equations,
and the evolution of the color charge.
Hence all dynamical equations arise from a single geometric object,
namely the symplectic form of the extended SHP phase space.
